\chapter{Arguing Requirements through the Lean Web Server}
\label{ch:web-server}

\index{vibecoding}
\index{requirements!durable}
The LeanFM web server presents a conversational workbench, but conversation is
not its product. Its product is a durable pair of artifacts: a typed Lean
requirements model and a protobuf description of the bytes exchanged by the
participants. The useful form of requirements vibecoding is therefore not
``ask a model to write the system.'' It is an argument in which a person keeps
challenging a proposed boundary until the remaining claims can be compiled,
checked, measured, or left as explicit proof obligations.

Chat history is evidence about how a decision was reached. It is not the
decision. A future reader must be able to discard the transcript, open the
generated files, and recover the actors, observable events, alternatives,
correlation rules, timing assumptions, resource bounds, and properties.

\section{The argument loop}

The interface resembles ChatGPT because natural language is convenient for
finding omissions. The loop, however, has a formal boundary:

\begin{center}
\begin{tikzpicture}[>=Latex,node distance=1.15cm,
  box/.style={draw,rounded corners,align=center,minimum width=3.1cm,
    minimum height=.75cm,fill=black!3},flow/.style={->,thick}]
  \node[box] (claim) {state or challenge\\a requirement};
  \node[box,right=of claim] (draft) {materialize a candidate\\Lean + protobuf pair};
  \node[box,below=of draft] (check) {compile, validate,\\simulate, and render};
  \node[box,left=of check] (evidence) {inspect traces, plots,\\violations, and obligations};
  \draw[flow] (claim) -- (draft);
  \draw[flow] (draft) -- (check);
  \draw[flow] (check) -- (evidence);
  \draw[flow] (evidence) -- node[left,align=right,font=\small]
    {counterexample or\\sharper question} (claim);
\end{tikzpicture}
\end{center}

The human contribution is often adversarial. ``The upload completes'' should
prompt: what observable event constitutes completion, who observes it, and does
it mean receipt, durable storage, or policy acceptance? ``The service handles
load'' should prompt: how many messages fit in each queue, how many
\texttt{(session, task)} pairs may be active, and what happens at the bound?
The assistant's job is to turn each answer into structure, not merely agreeable
prose.

A productive exchange might proceed as follows.

\begin{quote}\small
\textbf{Person:} A submitted upload eventually completes.\\
\textbf{Assistant:} I will treat an \texttt{uploadCompleted} response as
completion.\\
\textbf{Person:} No. Completion requires the storage actor's acknowledgment,
and the client must be able to distinguish rejection from timeout. Both queues
are finite.\\
\textbf{Assistant:} Then the grammar has accepted, rejected, and timed-out
alternatives; acceptance contains a storage acknowledgment before the response.
I will add queue capacities and a bounded-liveness obligation.\\
\textbf{Person:} Also show the behavior with many sessions in flight, estimate
the completion distribution from an event stream, and fail the requirement if
queued bytes exceed the memory budget.
\end{quote}

Every turn should change a named assumption, artifact, check, or open
obligation. If it changes none of them, it has not improved the specification.

\section{The server is a staging environment}

The executable \texttt{lake exe leanfm-server} serves the assistant workbench
on the loopback interface. An authenticated session has its own generated
\texttt{Requirements.lean}, \texttt{Requirements.proto}, and
\texttt{Implementation.lean}. The first two specify observable behavior and
values that resolve to bytes. They do not imply HTTP. The implementation file
selects how those bytes travel---for example HTTP method and path, an
in-process channel, TCP, or another adapter. The HTTP interface
can read and replace these files through
\texttt{/api/session/generated/requirements.lean} and
\texttt{/api/session/generated/requirements.proto}. Validation, rendered views,
metrics, reports, and OpenAPI descriptions are projections of the same model.

Session isolation is useful during an argument, but the current session files
live below \texttt{/tmp/leanfm-sessions}. That is a workspace, not durable
storage. A concluded argument must export or commit the accepted files---for
example as \texttt{LeanFM/LLMGenerated/Requirements.lean},
\texttt{LeanFM/LLMGenerated/Requirements.proto}, and the separate
\texttt{LeanFM/LLMGenerated/Implementation.lean}---together with any evidence
needed by the project. Durability means that the specification has an ordinary
name, history, review boundary, and reproducible checks independent of the chat
service.

The server may use a language model to propose edits. Compilation and validation
remain deterministic. The model is not the proof checker, and fluent prose is
not evidence that an obligation holds.

\section{What the Lean artifact must settle}

The Lean file is the semantic source. It records actors, messages, per-task
state machines, multiparty grammar productions, CSP-style processes, framing,
properties, proof obligations, probability summaries, charts, and explanatory
Markdown. A serious argument should force answers in each of the following
dimensions.

\begin{description}[style=nextline]
  \item[Observable behavior]
  Which actor sends and receives each event? Which fields can an observer see?
  A message carries its \texttt{session} and \texttt{task}, its clock timestamp,
  and a list of prior event identifiers. Thus a join may cite several prior
  events, while several events may fork from the same prior event.

  \item[Component interactions]
  Global grammar productions state legal conversations; actor projections state
  what each component may send or receive. Composition must preserve endpoints,
  message types, correlation, and causal references. A diagram is useful
  evidence, but the grammar is the compositional object.

  \item[Statistical properties]
  User choices make the input probabilistic even when every server component is
  deterministic. Those choices propagate through the composed system as an MDP.
  The specification distinguishes unconstrained alternatives from weighted
  choices, states dwell-time distributions, and names reducers for success,
  latency, throughput, and queue occupancy. An observed event stream can then
  estimate the same quantities and be compared with the model.

  \item[Queueing and concurrency]
  Every actor has finite inbound and outbound queues. Several
  \texttt{(session, task)} pairs may be in flight at once, so a requirement must
  not silently reason about only one conversation. Full and empty queues produce
  observable blocking or sleeping behavior, and transport is a separate step
  from production and receipt.

  \item[Memory pressure]
  Message-count capacity alone is insufficient when payload sizes vary. A
  checkable requirement should declare a byte budget and derive queued bytes from
  the encoded envelopes. If active tasks retain additional state, the model must
  declare a bounded footprint for that state too. A representative obligation is
  \[
    \mathrm{AG}\bigl(\operatorname{queuedBytes}(a) +
      \operatorname{inFlightBytes}(a) \leq \operatorname{budget}(a)\bigr).
  \]
  This is a modeling obligation for LeanFM, not a claim that the present
  validator already proves heap usage of an arbitrary implementation.
\end{description}

\section{Interrogate, characterize, synthesize, replay}

The conversational interface is not merely a code generator. Its first job is
to discover whether the desired claims are identifiable. It asks for scenario
and task boundaries, start-to-end backpointers, concrete actor populations,
clock and work units, observation windows, queue and memory limits, outage and
recovery events, and the provenance of probabilities. A missing answer is
reported as an indeterminate metric rather than replaced by an invented value.

Once those fields are available, one event stream should produce a standard
characterization out of the box:

\begin{itemize}
  \item one interaction diagram and derived state machine per scenario;
  \item XY line graphs for time-, load-, population-, and replica-indexed metrics;
  \item pie-chart histograms for categorical outcomes and distributions;
  \item uptime, outage fraction, reliability, and MTTR;
  \item client-experienced and server-aggregate throughput; and
  \item latency distributions and named percentiles.
\end{itemize}

Some streams contain observations; other requirements contain only expected
distributions. These are labeled separately. An expected distribution can drive
generation, but it is not presented as a measurement.

Generative similarity closes the loop. Reducers first turn the source stream
into a named metric vector, including histogram bins as individual exact ratios.
The generator produces a candidate scenario, serializes it to bytes, decodes it,
and runs the same reducers. The candidate is accepted only when every required
metric lies within its declared relative tolerance. Thus ``similar'' means
measurement-equivalent for the stated characterization, not visually plausible
or textually similar. Several very different event processes may satisfy the
same finite metric vector; the grammar, invariants, and byte-level language still
constrain which candidates are legal.

\section{The protobuf artifact is part of the requirement}

The parallel protobuf schema answers a question that a state diagram cannot:
what bytes can cross a component boundary? Atomic payload messages and a
\texttt{RequirementEnvelope.oneof} discriminator connect a grammar terminal to
its wire representation. Validation checks, among other things, that field tags
are nonzero and unique and that framing is explicit.

This pairing exposes useful disagreements. A Lean event may require a
\texttt{priorIds} list while the wire schema provides only one predecessor. A
latency property may depend on a clock that is never transmitted or recorded. A
memory budget may be stated in bytes while an unbounded protobuf field remains.
These are requirements defects, not serialization trivia.

The exact wire size is also an input to resource analysis:

\begin{lstlisting}
queuedBytes actor observation :=
  (observation.inbound actor).sum encodedSize
  + (observation.outbound actor).sum encodedSize

memoryPressure actor observation :=
  queuedBytes actor observation
  + inFlightFootprint actor observation.activeWork
\end{lstlisting}

The first term can be computed from concrete envelopes. The second requires a
declared abstraction of implementation state and must remain visibly labeled as
such.

\section{A ladder of checks}

No single green check answers every kind of requirement. The workbench should
present a ladder of evidence:

\begin{enumerate}
  \item Lean elaboration and protobuf parsing establish well-formed artifacts.
  \item Structural validation rejects missing actors, endpoints, terminals,
        framing, terminal states, process coverage, or real proof obligations.
  \item Grammar and projected state-machine checks cover ordering and component
        interaction.
  \item Temporal checking covers safety, reachability, and appropriately bounded
        liveness over observable states.
  \item Probability and dwell reducers predict completion probability and
        latency; event-stream reducers estimate them from observations.
  \item Queue and byte-budget checks cover capacity, backpressure, concurrency,
        and modeled memory pressure.
  \item Rendered interaction diagrams, state machines, plots, and reports make
        the same typed model reviewable by people.
\end{enumerate}

The distinction between model prediction and measurement must remain visible.
An incomplete observation window may censor slow tasks; a frequency estimate is
not automatically a proof of the declared probability; and a bounded trace does
not establish unbounded liveness. Good generated reports say which claim was
proved, calculated, sampled, or merely assumed.

\section{How to argue well}

The most useful prompts attack boundaries rather than ask for polish:

\begin{itemize}
  \item ``Show me the observable event that makes that sentence true.''
  \item ``Which alternative ends here, and how is it bracketed in the grammar?''
  \item ``What happens when this inbound or outbound queue is full?''
  \item ``Run two sessions and two tasks concurrently; which events correlate?''
  \item ``Which random choice belongs to the user, and how does it change the
        downstream completion and latency distributions?''
  \item ``Derive the plot from the event stream, including incomplete tasks.''
  \item ``Show the corresponding protobuf payload and its encoded-size bound.''
  \item ``Is this result proved, model-checked, simulated, measured, or assumed?''
\end{itemize}

The argument is finished enough to review when deleting the transcript loses no
normative information; a clean checkout can compile and validate the artifacts;
every diagram and plot is derived rather than independently drawn; failures have
observable outcomes; quantitative assumptions have units and scope; and every
unsettled claim is an explicit obligation. That is the requirements-centered
meaning of vibecoding in LeanFM: conversation discovers the specification, but
the checked artifacts outlive the conversation.
